\section{Artinian comparison and \'etale charts}\label{sec:comparison}

We now prove the locally quasi-finite criterion with the pro-modular premise. The first step explains exactly how Artinian tests determine a natural transformation once descent and effectivity are available.

\begin{lemma}\label{lem:comparison}
Let $G,F$ satisfy conditions \ref{ax:zariski}--\ref{ax:complete}, and let $G\to F$ be a natural transformation bijective on every local Artinian $S$-algebra. Then it is an isomorphism on all $S$-schemes.
\end{lemma}
\begin{proof}
First prove injectivity. Over a noetherian ring, two sections with equal image agree on every Artinian quotient at every maximal ideal. Condition~\ref{ax:complete} makes them agree on the corresponding complete local rings; fpqc injectivity then makes them agree on the local rings. Finitarity spreads those equalities to neighborhoods, and Zariski descent proves global equality. For an arbitrary algebra over a noetherian affine base, a pair of sections and equality of their images descend to a finite-type algebra by finitarity. This proves injectivity on arbitrary rings, including the nonnoetherian overlap rings that occur next.

For surjectivity over a noetherian local $A$, an element of $F(A)$ has unique compatible preimages over $A/\mathfrak m^n$. They give an element of $G(\widehat A)$. Its two images over $\widehat A\otimes_A\widehat A$ agree by the injectivity just proved. Fpqc descent gives an element of $G(A)$ with the prescribed image. Localizing and spreading extends the conclusion to noetherian rings, and finite-type approximation then extends it to all rings. Zariski gluing finishes the proof.
\end{proof}

A scheme locally of finite presentation satisfies the effectivity condition used here: a compatible sequence of maps from the Artinian quotients of a complete local ring lands in a fixed affine neighborhood of their common closed point. Taking the inverse limit of the coordinate homomorphisms gives the map from the complete local ring. In particular no properness hypothesis is needed for this comparison.

Condition~\ref{ax:artin}, combined with Lemma~\ref{lem:comparison}, has an important consequence: over a local Artinian base, $F$ is represented on \emph{all} test schemes by a disjoint union of finite local schemes. The finite representing algebras in condition~\ref{ax:artin} are themselves Artinian and carry their universal elements, so they give the required natural transformation before the comparison lemma is applied.

\begin{lemma}[Artinian form of the \'etale criterion]\label{lem:etale}
Over a locally noetherian base, conditions \ref{ax:zariski}--\ref{ax:complete} and \ref{ax:valuation}, together with unique lifting through nilpotent ideals on local Artinian tests, imply representability by a separated \'etale scheme.
\end{lemma}
\begin{proof}
Murre's Corollary 2 \cite[p.~247]{Murre} proves this with unique lifting on all nilpotent tests. We verify the reduction, beginning with uniqueness. For a noetherian local pair $(A,I)$, equality modulo the nilpotent ideal $I$ implies equality modulo every $\mathfrak m^n$: apply Artinian uniqueness to $A/\mathfrak m^n\to A/(I+\mathfrak m^n)$. Condition~\ref{ax:complete} and faithful flatness of completion give equality over $A$. Localization gives uniqueness for noetherian rings.

For an arbitrary algebra $A$ over a noetherian affine part of the base, write
\[
 A/I=\varinjlim_{\substack{J\subseteq I\\J\text{ finitely generated}}} A/J.
\]
An equality of two sections over $A/I$ already holds over one $A/J$ by finitarity. Choose generators of $J$ and an exponent $n$ with $I^n=0$. The two sections, their equality modulo $J$, these generators and the finitely many equations making all degree-$n$ monomials in them zero descend to a finite-type algebra. Its corresponding ideal is nilpotent. Uniqueness there proves uniqueness over $A$. Thus uniqueness holds also on the possibly nonnoetherian completion overlaps.

For existence over a noetherian local $A$, restrict a section over $A/I$ to $A/(I+\mathfrak m^n)$ and lift it uniquely to $A/\mathfrak m^n$. The lifts are compatible and give a section over $\widehat A$. Its reduction agrees with the original section over $\widehat A/I\widehat A$, by effectivity for that complete noetherian quotient. Its two pullbacks to $\widehat A\otimes_A\widehat A$ have equal reductions modulo $I$; the uniqueness already proved makes them equal. Fpqc descent gives the lift over $A$. Localization, spreading and Zariski gluing handle noetherian rings. Finally, a section over an arbitrary $A/I$ comes from one $A/J$ in the displayed limit. Descend that section, generators of $J$ and their nilpotence equations to a finite-type stage and apply the noetherian existence assertion. Pullback gives the required lift over $A$. Murre's corollary now applies; its remaining hypotheses are exactly the stated sheaf, finitary, complete-local and valuative conditions.
\end{proof}

For a finite-type pair $(T,\eta)$, let $E(T,\eta)$ be its \emph{Artinian formally \'etale locus}: at a point of $E$, every diagram of a nilpotent closed immersion of local Artinian schemes centered at that point, a map into $T$ on the closed subscheme, and an element of $F$ extending its image has a unique lifting map to $T$. Residue-field extensions are allowed. This property is invariant under \'etale changes of source and base and can be checked after such a change: lift the residue point after a finite \'etale field extension, apply formal \'etaleness of the change, and descend the unique map. Write $M(T,\eta)$ for the local modularity locus of Definition~\ref{def:modular}.

\section{Openness of the modularity locus}\label{sec:openness}

This section proves openness with propagation relative to complete local bases. The main issue is to control infinitesimal directions transverse to a closed stratum; a statement only about its reduced generic point does not do that.

\begin{lemma}\label{lem:uniform}
Let $R$ be noetherian, $I\subset R$ an ideal, and $C$ a finite $R$-module. There is $N\geq1$ such that the following holds. If $R\to R'$ is flat, $R'$ is noetherian local and $IR'$ is contained in its maximal ideal, $D$ is finite over $R'$, and $\phi:D\to C\otimes_RR'$ is an isomorphism modulo $I^N$, then $\phi$ is an isomorphism. The same $N$ works for any direct-summand target of $C\otimes_RR'$.
\end{lemma}
\begin{proof}
Choose $0\to K\to R^d\to C\to0$. Artin--Rees gives $N$ with $K\cap I^NR^d\subset IK$ \cite[Tag 00IN]{Stacks}. Flat base change preserves the presentation and this intersection. Nakayama makes $\phi$ surjective. Lift the chosen generators to $D$; they generate modulo $I^N$ and hence generate $D$. If $K'$ is the kernel of $R'^d\to D$, then $K'\subset K\otimes_RR'$ and the isomorphism modulo $I^N$ gives
\[
 K\otimes_RR'\subset K'+I^NR'^d.
\]
Intersecting with the kernel and using Artin--Rees gives $K\otimes_RR'\subset K'+I(K\otimes_RR')$. Nakayama on the finite quotient proves equality of kernels. For a direct-summand target, add its complementary summand to the source and apply the result already proved.
\end{proof}

\begin{lemma}\label{lem:generization}
Assume conditions \ref{ax:zariski}--\ref{ax:valuation} and $(7_{\mathrm{rel}})$. Every generization of a point of $M(T,\eta)$ belongs to $E(T,\eta)$.
\end{lemma}
\begin{proof}
Let $x\in M(T,\eta)$ map to $s\in S$. Over $\kappa(s)$ the functor is represented by finite local schemes, by the Artinian comparison. After base change, the induced pair is locally modular for the restricted functor: use the original local universal property and pair its unique map with the new base map; the minimal-field condition identifies the required point of the fibre product. Therefore $T_s$ has a local-ring isomorphism to a finite representing branch at $x$, and $T/S$ is quasi-finite there.

Put $A=\OO_{S,s}$ and pass to $\widehat A$. The branch of $T_{\widehat A}$ at a point over $x$ has a finite local open neighborhood $C=\Spec B$, by the henselian finite-branch theorem \cite[Tag 04GH]{Stacks}. Its closed point is locally modular for $F_{\widehat A}$ on $\widehat A$-schemes. Condition $(7_{\mathrm{rel}})$, applied with complete base $\widehat A$, makes it so everywhere for that same restricted functor. For any noetherian $\widehat A$-test $(Y,\zeta)$ the fibre product $h_C\times_{F_{\widehat A}}h_Y$ satisfies the sheaf, finitary and effectivity conditions. Its Artinian lifting is unique by modularity, and valuative uniqueness follows from separation of the finite scheme $C$ over the base. Lemma~\ref{lem:etale} represents it by a separated \'etale $Y$-scheme. Finitary approximation handles arbitrary affine tests, and gluing handles all tests. Hence $h_C\to F_{\widehat A}$ is representable and \'etale.

Let $X=\Spec\OO_{T,x}$. Every function invertible at $x$ is a unit on the local scheme $C$, so $C$ is an open subscheme of $X\times_A\widehat A$. The map $C\to X$ is flat local on rings and therefore fpqc. Unique formal lifting for $C$ descends to $X$ as follows. For a nilpotent immersion $Z_0\subset Z$ and a map $Z_0\to X$, the inverse image of $C$ is an open $V_0$ of $(Z_0)_{\widehat A}$. It extends uniquely to an open $V$ of $Z_{\widehat A}$, and $V\to Z$ is fpqc: flatness holds because $V$ is open in the flat base change $Z\times_A\widehat A$, while surjectivity and quasi-compactness follow from the corresponding properties after nilpotent reduction. Lift over $V$ to $C$. On the overlap, pair both resulting maps to $X$ with the base structure from the first factor. Their reductions land in the same open $C$, so the maps do as well. Formal uniqueness for $C$ makes them equal. Fpqc descent then gives the unique map over $Z$.

Thus $h_X\to F_A$ is formally \'etale. A local Artinian test centered at a generization of $x$, and any nilpotent lifting of it into $T$, factor through $X$. This gives exactly the required Artinian lifting property at that generization.
\end{proof}

\begin{lemma}\label{lem:stratum}
Assume conditions \ref{ax:zariski}--\ref{ax:valuation} and \ref{ax:generic}. If an irreducible closed subset $Z\subset T$ has generic point $z\in E(T,\eta)$, then a nonempty relatively open subset of $Z$ is contained in $E(T,\eta)$.
\end{lemma}
\begin{proof}
Over the residue field of the image $s$ of $z$, the functor is represented by finite local schemes. The infinitesimal criterion makes the fibre map \'etale at $z$, so $T/S$ is quasi-finite there. After an \'etale neighborhood of $s$ and restriction around a point over $z$, we may take $T=\Spec C$ finite over an affine base $\Spec R$. To justify the replacement, isolate the finite local branch over the strict henselization using Tag 04GH; its finite presentation and its open inclusion descend to an \'etale neighborhood of finite presentation. At the same time the finite separable residue extension of the \'etale fibre map can be split. Over a separably closed base residue field all finite local branches have purely inseparable residue fields \cite[Tag 06DD]{Stacks}, so an \'etale map between them has equal residue fields. The required finite residue splitting descends to an \'etale neighborhood. These changes have open images and preserve and reflect the lifting condition.

Let $I\subset R$ define the closure of the image of $z$, and choose $N$ from Lemma~\ref{lem:uniform} for the finite module $C$. Over the Artinian local ring $(R/I^N)_I$, the functor is represented. The branch at $z$ maps \'etale to its finite local representing branch with identical residue field; hence that local map is an isomorphism. Remove the other irreducible components from $T\times_R\Spec(R/I^N)$. On the resulting irreducible open the generic point is locally modular over the irreducible noetherian base $R/I^N$. Condition~\ref{ax:generic} supplies a nonempty open $W$ on which it is locally modular.

Fix $w\in W$ with base image $q$, and put $R'=\widehat{R_q}$. The finite algebra $C\otimes_RR'$ splits into finite local factors; let $C_w$ be the factor through $w$. Conditions \ref{ax:complete} and \ref{ax:artin} give a finite complete local $R'$-algebra $D_w$ representing the local Artinian deformation functor at the selected field point. Here finiteness of the inverse limit follows by lifting a finite set of module generators at level one and correcting their coefficients successively in powers of the maximal ideal. Completeness gives a surjection from a finite free $R'$-module, and identifies its reductions with the finite-level representing algebras. The element on $C_w$ induces a homomorphism $D_w\to C_w$.

Local modularity over $R/I^N$ identifies these two complete local deformation algebras modulo $I^N$. Lemma~\ref{lem:uniform}, including its direct-summand assertion for the target $C_w$, makes the homomorphism itself an isomorphism. Hence the original pair is Artinian formally \'etale at $w$. Indeed centered Artinian tests factor uniquely through the corresponding completions, so this conclusion is precisely the definition of $E$.

If the \'etale alternative to the conclusion of condition~\ref{ax:generic} is used, the residue extension of $D_w/I^ND_w\to C_w/I^NC_w$ is finite separable. Take the finite \'etale local extension $D'_w$ of the henselian ring $D_w$ having that residue field. Its map to $C_w$ lifts by henselian uniqueness. Modulo $I^N$ it is an isomorphism, by the unique Artinian lifting properties and completeness. Lemma~\ref{lem:uniform} makes $D'_w\to C_w$ an isomorphism. Thus $D_w\to C_w$ is \'etale, which gives the same conclusion.

For a finite morphism, primes over $I$ are incomparable. Thus near $z$ the chosen irreducible component of the inverse image of $V(I)$ is exactly the transformed stratum $Z$. The open $W$ supplies the required relative open after the modifications, and its \'etale image gives a nonempty relative open in the original $Z$.
\end{proof}


\begin{proof}[Proof of Theorem~\ref{thm:openness}]
Work on a noetherian affine piece of $T$ over an affine piece of $S$, and let $U$ be the interior of $E=E(T,\eta)$. We first prove $M(T,\eta)\subset U$. If a modular point $x$ lay outside $U$, choose an irreducible component $Z$ of the closed complement containing it, with generic point $z$. Lemma~\ref{lem:generization} puts $z$ in $E$. Lemma~\ref{lem:stratum} gives a nonempty relative open $W\subset Z$ lying in $E$. Delete the other components of the complement and shrink around $z$ until its intersection with the complement lies in $W$. This open neighborhood lies entirely in $E$: its points lie either in $U$ or in $W$. By definition of the interior it lies in $U$, a contradiction.

On $U$, each fibre product $h_U\times_Fh_Y$ is represented and \'etale by Lemma~\ref{lem:etale}, first for noetherian $Y$, then for arbitrary tests by finitarity. Its image $Q$ is an open subfunctor of $F$, and the \'etale equivalence relation $U\times_FU\rightrightarrows U$ presents an algebraic space $Q$. Every point of $U$ is quasi-finite over $S$, as is seen on the represented Artinian residue fibre. Thus $Q/S$ is locally quasi-finite.

We verify separation before using the scheme criterion. For affine chart pieces $U_i,U_j$, the relation $U_i\times_QU_j\to U_i$ is separated and \'etale, and its fibres are finite because their points inject into a fibre of the finite-type quasi-finite scheme $U_j/S$. A separated \'etale map with finite fibres over a noetherian scheme is quasi-compact. In fact the loci where its fibre degree is at least $d$ are open \cite[Tag 086R]{Stacks}; the finite degrees at the finitely many generic points bound all degrees. Tag 07RY then gives a finite locally closed stratification on which the map is finite locally free. Each inverse image is quasi-compact, proving the claim. Hence the relation is quasi-compact and $Q/S$ is quasi-separated. Valuative uniqueness makes it separated \cite[Tag 03KV]{Stacks}. A separated locally quasi-finite algebraic space over a scheme is a scheme \cite[Tag 0417, Proposition 67.50.2]{Stacks}.

On $U$, modularity is now the local-ring-isomorphism locus of $U\to Q$. This is open by finite presentation and spreading of the inverse, as in Proposition~\ref{prop:finite-type}. Since every modular point lies in $U$, this proves the assertion on the affine piece, hence globally. In the \'etale alternative to $(7_{\mathrm{rel}})$, Lemma~\ref{lem:generization} already has exactly the \'etale chart it needs; the preceding proof is unchanged.
\end{proof}

\section{The locally quasi-finite analogue}\label{sec:criterion}



The finiteness in condition~\ref{ax:artin} allows formal branches to be assembled as finite algebras. The strengthened propagation premise makes their restrictions to the punctured base \'etale. Valuative uniqueness prevents different branches from meeting there, and the openness theorem turns local representations into global charts. This explains where each hypothesis enters the proof.

\begin{lemma}\label{lem:local-global}
Suppose conditions \ref{ax:zariski}--\ref{ax:valuation}, $(7_{\mathrm{rel}})$, and \ref{ax:generic} hold and $F$ is represented by a separated locally quasi-finite scheme over each $\Spec\OO_{S,s}$. Then it is represented by such a scheme over $S$.
\end{lemma}
\begin{proof}
Take an affine finite-type open of a local representing scheme. Its finite presentation and universal element descend, by finitarity, to an affine finite-type pair $(V,\xi)$ over a neighborhood of $s$. Over the local base it is modular at every point. These points are also modular over $S$: a local $S$-test centered at a generization of $s$ factors uniquely through $\Spec\OO_{S,s}$. Theorem~\ref{thm:openness} makes its modularity locus open in $V$. Its closed complement has empty localization at $s$; finite presentation and quasi-compactness allow one denominator outside $s$ to remove it. After shrinking the base, $h_V\to F$ is modular everywhere and in particular representable \'etale, by Lemma~\ref{lem:etale} applied to its fibre products.

Taking these charts for all local points gives a surjective \'etale presentation of $F$ and hence an algebraic space. It is locally quasi-finite. The relation between two affine chart pieces has finite fibres over either piece, checked on the local representing schemes. The finite-fibre argument in Theorem~\ref{thm:openness} makes the relation quasi-compact. Valuative uniqueness gives separation, and Tag 0417 gives a scheme. This proves the lemma.
\end{proof}

\begin{proof}[Proof of Theorem~\ref{thm:criterion}]
We prove sufficiency by induction on the dimension of noetherian bases, simultaneously after their noetherian base changes. Over an Artinian base, condition~\ref{ax:artin} and Lemma~\ref{lem:comparison} give the representation. Suppose the assertion is known in dimensions less than $n$, and first take a complete noetherian local base $(A,\mathfrak m)$ of dimension $n$. Its punctured spectrum $U$ has dimension at most $n-1$, so $F_U$ is represented by a separated locally quasi-finite $X_U$.

For each component of the represented closed fibre, condition~\ref{ax:artin} gives compatible finite local algebras $B_{i,r}$ over $A/\mathfrak m^r$. Their limit $B_i$ is finite local over $A$: lift generators from the first level and correct coefficients successively in powers of $\mathfrak m$ as in Lemma~\ref{lem:stratum}. This also proves $B_i/\mathfrak m^rB_i=B_{i,r}$. Condition~\ref{ax:complete}, applied to the finite complete local ring $B_i$, supplies the universal element $\xi_i$. The two adic filtrations on this ring are cofinal. Condition $(7_{\mathrm{pro}})$ makes the maps $Z_i=\Spec B_i\to F$ \'etale, in particular over $U$.

The maps $(Z_i)_U\to X_U$ are open immersions with pairwise disjoint images. Indeed, an off-diagonal point of their equality relation has closure in the finite $A$-scheme $Z_i\times_AZ_j$ meeting the closed fibre. A noetherian valuative specialization gives a trait from a point of that relation to a point in the closed fibre, after allowing a residue extension. The two sections of $F$ agree at its generic point and hence everywhere by condition~\ref{ax:valuation}. At the closed point their Artinian universal properties force the same closed-fibre component, and for $i=j$ force the two maps $B_i\to V/\mathfrak m_V^r$ to agree for every $r$. Krull intersection in the discrete valuation ring $V$ makes their maps to $V$ equal, contradicting the chosen off-diagonal generic point. Each map $(Z_i)_U\to X_U$ is separated: its source is separated over $U$, since $Z_i$ is finite over $A$, and $X_U$ is an $U$-scheme. Thus its diagonal, already open by \'etaleness, is also closed. Hence the equality relations have no complement to the diagonal as schemes. The maps are \'etale monomorphisms, thus open immersions, and different images are disjoint.

Glue $\coprod_iZ_i$ to $X_U$ along these open immersions. The resulting scheme $X$ is locally of finite type and carries a section of $F$. Its natural map $h_X\to F$ is bijective on local Artinian tests: tests centered at the closed point use the $B_{i,r}$, and the others factor through $U$. Lemma~\ref{lem:comparison} makes it an isomorphism on all tests. It is locally quasi-finite. The gluing has quasi-compact intersections on noetherian affine pieces, since distinct finite branches are disjoint over $U$; thus it is quasi-separated. Valuative uniqueness gives separation.

For a noetherian local base that is not complete, its faithfully flat completion gives this representation and a descent datum. Murre's Appendix, Theorem, case (d) \cite[\S4, pp.~259--260]{Murre}, makes that datum effective: the cover is fpqc, its source is noetherian, and the represented scheme is separated, locally quasi-finite and locally of finite presentation. Equivalently, a separated locally quasi-finite morphism is ind-quasi-affine, and such morphisms have effective fpqc descent \cite[Tags 0AP9 and 0APK]{Stacks}. These properties descend. Thus all local representations in dimension $n$ exist. Condition $(7_{\mathrm{pro}})$ implies $(7_{\mathrm{rel}})$, so Lemma~\ref{lem:local-global} gives the global assertion and completes the induction. Every noetherian local ring has finite dimension. Applying the same local-to-global lemma therefore also treats locally noetherian bases whose dimensions are not globally bounded.

Conversely let $F=h_X$ with $X/S$ separated and locally quasi-finite. The sheaf, finitary, complete-local, Artinian and valuative conditions are the corresponding standard properties of this morphism. For $(7_{\mathrm{pro}})$, over a complete local base the branch of $X$ at a closed-fibre point is finite local and complete by the henselian finite-branch theorem. A finite algebra with the same Artinian universal property has the same quotients at every level, so completeness identifies it with that branch. Its map to $X$ is an open immersion, hence locally modular everywhere. Condition~\ref{ax:generic} follows from finite presentation and spreading a local-ring isomorphism. This proves necessity.
\end{proof}

The argument realizes Murre's announced locally quasi-finite analogue with a specified pro-modular propagation axiom. It uses finite algebras, fpqc completion descent and a uniform Artin--Rees test. It does not invoke Artin approximation over a nonexcellent base.
